%%%%%%%%%%%%%%%%%%%%%%%%%%%%%%%%%%%%%%%%
% CAPÍTULO - INTRODUÇÃO

\chapter{INTRODUÇÃO}
\label{chp:intr}


A viagem e a exploração espacial incluem o transporte de veículos através do espaço sideral, e estão associados às tecnologias de propulsão, astrodinâmica, astronáutica, e construção e lançamento de veículos. O interesse pelos astros e sistema solar remonta o início da história da humanidade, e as pessoas dependiam de medições indiretas e observações da luz visível e à distância dos objetos. Com a evolução da tecnologia e do conhecimento, foi possível o envio de espaçonaves ao espaço com diversas finalidades, incluindo exploração, observação, telecomunicação, defesa, e outros.



\section{Objetivo}

O objetivo deste trabalho é explorar duas missões espaciais, NanosatC-Br1 e OSIRIS-REx, mostrando as missões, descrição, funcionamento, e equipamentos da espaçonave.



\section{Organização do texto}

A seguir está uma uma breve descrição dos demais capítulos deste documento, mostrando
a sua organização:



\begin{itemize}


\item Capítulo \ref{chp:intr}: Introdução.

\item Capítulo \ref{chp:nsat}: NanosatC-Br1.

\item Capítulo \ref{chp:orex}: OSIRIS-REx.

\item Capítulo \ref{chp:conc}: Comentários finais.


\end{itemize}



%% EOF